% Part: lambda-calculus
% Chapter: syntax
% Section: unique-readability

\documentclass[../../../include/open-logic-section]{subfiles}

\begin{document}

\olfileid{lam}{syn}{unq}
\olsection{وحدانية القراءة}

هل يمكن أن تتعدد طرق تكوين الحد الواحد؟ سنثبت أنها لا تتعدد:
فلكل حد لامبدا وجه واحد في البناء والتأويل. ومرادنا بـ\emph{التكوين}
فيما يلي بناء الحد بتطبيق قواعد \olref[trm]{defn:term} مرة أو مرات.

\begin{lem}\ollabel{lem:term-start}
لا يخلو أول الحد من أن يكون متغيرًا أو قوسًا.
\end{lem}

\begin{proof}
إنما يكون الشيء حدًا إذا بني بقواعد~\olref[trm]{defn:term}.
فما تبنيه \olref[trm]{defn:term-var} متغير، وما تبنيه
\olref[trm]{defn:term-abs} أو~\olref[trm]{defn:term-app} أوله قوس.
\end{proof}

\begin{lem}\ollabel{lem:app-start}
أول ناتج التطبيق قوسان، أو قوس يتلوه متغير.
\end{lem}

\begin{proof}
ليكن~${\OLId{latin-looped}{M}}$ حاصل تطبيق؛ فصورته~$(PQ)$، وأوله إذن قوس.
ثم إن~${\OLId{latin-looped}{P}}$ حد، فأوله، بموجب~\olref{lem:term-start}، قوس أو متغير.
\end{proof}

\begin{lem}\ollabel{lem:initial}
كل جزء ابتدائي حقيقي من حد يمتنع أن يكون حدًا.
\end{lem}

\begin{prob}
أقم برهان~\olref[lam][syn][unq]{lem:initial} مستعملًا الاستقراء على طول الحد.
\end{prob}

\begin{prop}[وحدانية القراءة] \ollabel{prop:unq}
تكوين كل حد وحيد؛ فإن انتهى تكوين إلى حد~${\OLId{latin-looped}{M}}$، امتنع أن ينتهي إليه تكوين آخر.
\end{prop}

\begin{proof}
  نستقرئ تكوين الحدود، ونميز الحالات الثلاث:
  \begin{enumerate}
    \item إذا كان~${\OLId{latin-looped}{M}}$ على صورة~${\OLId{latin-ordinary}{x}}$، حيث~${\OLId{latin-ordinary}{x}}$ متغير،
      لم يدخل التجريد ولا التطبيق في بنائه؛ إذ يبدأ ناتج كل منهما بقوس.
      فلا سبيل إلى بناء~${\OLId{latin-looped}{M}}$ إلا أن يكون تكوين~${\OLId{latin-looped}{M}}$ خطوة واحدة من
      \olref[trm]{defn:term}\olref[trm]{defn:term-var}.
    \item إذا كان~${\OLId{latin-looped}{M}}$ على صورة~$(\lambd[{\OLId{latin-ordinary}{x}}][{\OLId{latin-looped}{N}}])$، حيث~${\OLId{latin-ordinary}{x}}$ متغير و${\OLId{latin-looped}{N}}$ حد،
      امتنع أن تبنيه القاعدة
      \olref[trm]{defn:term}\olref[trm]{defn:term-var}؛ لأنه ليس متغيرًا مفردًا.
      ويمتنع أيضًا أن يكون حاصل تطبيق، بموجب~\olref{lem:app-start}.
      فليس~${\OLId{latin-looped}{M}}$ إلا حاصل تجريد~${\OLId{latin-looped}{N}}$، وتكوين~${\OLId{latin-looped}{N}}$ وحيد بفرض الاستقراء.
    \item إذا كان~${\OLId{latin-looped}{M}}$ على صورة~$(PQ)$، حيث~${\OLId{latin-looped}{P}}$ و${\OLId{latin-looped}{Q}}$ حدان،
      امتنع بناؤه بالقاعدة
      \olref[trm]{defn:term}\olref[trm]{defn:term-var}، لأن أوله قوس.
      وتقتضي~\olref{lem:term-start} أن~${\OLId{latin-looped}{P}}$ لا يبدأ بـ$\lambd$؛
      فلا يكون~$(PQ)$ حاصل تجريد. فلنفرض أن~${\OLId{latin-looped}{M}}$ بني بتطبيق آخر،
      فكانت صورته أيضًا~$({\OLId{latin-looped}{P}}'{\OLId{latin-looped}{Q}}')$. إذا كان~${\OLId{latin-looped}{P}} \neq
      {\OLId{latin-looped}{P}}'$، لزم أن يكون أحد~${\OLId{latin-looped}{P}}$ و${\OLId{latin-looped}{P}}'$
      جزءًا ابتدائيًا حقيقيًا من الآخر، وهو ممتنع بموجب~\olref{lem:initial}.
      إذن~${\OLId{latin-looped}{P}}={\OLId{latin-looped}{P}}'$، ومنه~${\OLId{latin-looped}{Q}}={\OLId{latin-looped}{Q}}'$. وبذلك يتعين~${\OLId{latin-looped}{P}}$ و${\OLId{latin-looped}{Q}}$،
      ثم ينفرد كل من~${\OLId{latin-looped}{P}}$ و${\OLId{latin-looped}{Q}}$ بتكوينه، بمقتضى فرض الاستقراء.
  \end{enumerate}
\end{proof}

وللقضية نفسها عبارة أيسر في القراءة:
\begin{prop}
  ينحصر الحد~${\OLId{latin-looped}{M}}$ في إحدى الصور التالية:
  \begin{enumerate}
    \item ${\OLId{latin-ordinary}{x}}$، وفيها~${\OLId{latin-ordinary}{x}}$ متغير يتعين من~${\OLId{latin-looped}{M}}$ على وجه واحد.
    \item $(\lambd[{\OLId{latin-ordinary}{x}}][{\OLId{latin-looped}{N}}])$، وفيها~${\OLId{latin-ordinary}{x}}$ متغير و${\OLId{latin-looped}{N}}$ حد آخر، يتعين كلاهما من
      ${\OLId{latin-looped}{M}}$ على وجه واحد.
    \item $(PQ)$، وفيها~${\OLId{latin-looped}{P}}$ و${\OLId{latin-looped}{Q}}$ حدان يتعينان من~${\OLId{latin-looped}{M}}$ على وجه واحد.
  \end{enumerate}
\end{prop}

\end{document}

